% !PRINTED 330
\weberpage{330}

\begin{center}
\gesperrt{Zehnter Abschnitt.}\\[4pt]
\textbf{Algebraische Zahlen und Formen.}
\end{center}

\begin{center}
\rule{0.25\textwidth}{0.4pt}
\end{center}

\begin{center}
\S\,87. \gesperrt{Ideale und Formen in algebraischen Körpern.}
\end{center}

Das Interesse, das die Theorie der elliptischen Funktionen für den Algebraiker hat, entspringt aus ihrer Anwendung auf die Theorie der quadratischen Irrationalzahlen, zu denen sie in einer analogen Beziehung stehen, wie die Einheitswurzeln zu den rationalen Zahlen. Sie bieten uns das erste und bisher einzige in seinen Gesetzen genauer bekannte Beispiel eines Gebietes algebraischer Zahlen, die über die Kreisteilungszahlen hinausgehen. Um die Theorie dieser Zahlen eingehender darstellen zu können, müssen wir einige Sätze aus der Theorie der Formen und algebraischen Körper und speziell der quadratischen Formen vorausschicken.

In Bd. II, \S\,163, 169 hat sich ein Zusammenhang ergeben zwischen den Idealen eines algebraischen Körpers $\Omega$ vom $n$ten Grade und gewissen homogenen Funktionen $n$ten Grades von $n$ Variablen. Ehe wir dies auf quadratische Körper anwenden, sei kurz an die allgemeinen Sätze erinnert.

Es sei $\mathfrak{a}$ ein Ideal eines Körpers $\Omega$ und $\alpha_1, \alpha_2, \ldots, \alpha_n$ eine Basis dieses Ideals. Ferner sei $\omega_1, \omega_2, \ldots, \omega_n$ eine Minimalbasis von $\Omega$. Dann ist
\weq{1}{(\alpha_1, \alpha_2, \ldots, \alpha_n) = (A)(\omega_1, \omega_2, \ldots, \omega_n),}
wenn $(A)$ eine \gesperrt{lineare Substitution mit ganzzahligen Koeffizienten} bedeutet. Um zu einer anderen Basis $\alpha_1', \alpha_2', \ldots, \alpha_n'$ von $\mathfrak{a}$ überzugehen, mache man eine lineare Substitution $L$ mit der Determinante $\pm 1$,
\weq{2}{(\alpha_1', \alpha_2', \ldots, \alpha_n') = (L)(\alpha_1, \alpha_2, \ldots, \alpha_n),}
woraus sich ergibt
\weq{3}{(\alpha_1', \alpha_2', \ldots, \alpha_n') = (L)(A)(\omega_1, \omega, \ldots, \omega_n)^{\text{1)}}.} % [sic] printed \textquotedbl{}$\omega$\textquotedbl{} without subscript 2 in the second entry

\medskip
\noindent\rule{0.3\textwidth}{0.4pt}\\
$^{1}$) Über lineare Substitutionen und ihre Zusammensetzung vergleiche man den sechsten Abschnitt des zweiten Bandes.
